% Generated draft: requires mathematical and source review.
\documentclass[11pt]{article}
\usepackage[margin=25mm]{geometry}
\usepackage{fontspec}
\setmainfont{Latin Modern Roman}
\usepackage{amsmath,amssymb,mathtools,unicode-math}
\setmathfont{Latin Modern Math}
\usepackage{enumitem,needspace,xurl,hyperref}
\setlength{\parindent}{0pt}
\setlength{\parskip}{5pt}
\setlength{\emergencystretch}{3em}
\newcommand{\N}{\mathbb{N}}
\newcommand{\Z}{\mathbb{Z}}
\newcommand{\Q}{\mathbb{Q}}
\newcommand{\R}{\mathbb{R}}
\newcommand{\C}{\mathbb{C}}
\newcommand{\F}{\mathbb{F}}
\newcommand{\A}{\mathbb{A}}
\newcommand{\PP}{\mathbb{P}}
\newcommand{\E}{\mathbb{E}}
\newcommand{\eps}{\varepsilon}
\newcommand{\dd}{\,\mathrm{d}}
\begin{document}
\section*{UnsolvedMath ID: 30002449}
\textbf{Source problem code:} catalogue entry
\subsection{Definitions and mathematical statement}
For an ordinary commutative interpretation, a cluster algebra is an algebra generated over a specified coefficient ring by the cluster variables obtained from an initial seed by iterated seed mutation. Once the coefficient ring and cluster algebra conventions are fixed, a cluster algebra structure on an algebra \(R\) includes an algebra isomorphism \(R\cong\mathcal{A}\), where \(\mathcal{A}\) is a cluster algebra with those conventions. This interpretation supplies terminology only; the source does not specify these choices or exclude other variants.

The supplied question is: \emph{Which algebras carry natural cluster algebra structures?}

The intended task is to characterize the algebras admitting such structures. However, the source specifies neither the class of algebras under consideration nor the meaning of ``natural.'' It also leaves the base ring or field, coefficient conventions, and treatment of frozen variables unspecified. Consequently, existence of an abstract isomorphism to a cluster algebra cannot by itself be identified with the requested naturality condition. A precise classification statement requires these choices and a definition of naturality; they remain unresolved in this draft.
\subsection{Short English statement}
Determine which algebras admit cluster algebra structures that are natural in the intended setting. The supplied question leaves both the class of algebras and the naturality requirement undefined, so it does not yet specify a precise classification problem.
\subsection{Sources}
\begin{itemize}
\item Supplied local UnsolvedMath registry record, permanent ID \(30002449\), title \emph{Algebras with Natural Cluster Structures}, source problem code \texttt{catalogue entry}, in the Oberwolfach Reports Open Problems collection. No individual proposer is identified; the supplied proposed year is 2013.
\item Supplied report citation: \emph{Cluster Algebras and Related Topics} (2014), Oberwolfach Reports, pp.\ 3379--3432, DOI \texttt{10.4171/owr/2013/58}. Supplied source URL and literature reference labeled ``original OWR report'': \url{https://doi.org/10.4171/owr/2013/58}. The background describes the reports as documenting mathematical workshops and their open-problem discussions. The citation has not been independently consulted.
\item Supplied literature assessment, checked \texttt{2026-08-21}: the record labels the problem open and states that natural cluster structures are known for many important algebras, while the requested general classification is incomplete. This is supplied dated curation, not an independently verified assessment of mathematical or literature status; the record explicitly requests independent verification of time-sensitive status and cited claims.
\end{itemize}
\end{document}
